% 第三章 系统需求分析

\chapter{系统需求分析}


\section{应用场景分析}

FreeFlow面向链上音乐资源的发行和消费场景。创作者可以在桌面端完成音乐信息编辑、素材上传、metadata生成、链上发布和收益分账配置；听众可以检索链上音乐资源，查看作品详情，购买访问权并进行播放、评论和下载。系统服务端负责连接桌面端、IPFS存储和区块链合约，提供身份认证、资源索引和业务数据管理能力。


\section{用户角色分析}

\begin{table}[htbp]
    \linespread{1.5}
    \zihao{5}
    \songti
    \centering
    \caption{用户角色与需求说明}\label{tab:user_roles}
    \begin{tabular}{|l|p{10cm}|}
        \hline
        角色   & 主要需求                                                          \\ \hline
        创作者  & 创建发行草稿、编辑音乐信息、上传音频和封面、生成metadata、设置访问模式、设置收益分账、发布作品到链上、查看发布状态 \\ \hline
        听众   & 搜索链上音乐、查看作品详情、查询访问状态、购买访问权、播放音乐、加入链上音乐库、下载音乐、发表评论             \\ \hline
        系统服务 & 管理SIWE会话、代理Pinata上传、保存资源索引、记录购买和评论、提供检索与排序接口                  \\ \hline
    \end{tabular}
\end{table}


\section{功能需求分析}

系统功能需求如表\ref{tab:functional_requirements}所示。本文采用“需求分析 + 总体设计 + 核心模块实现 + 系统测试”的组织方式展开说明，各功能模块的具体实现与验证将在后续章节中进一步给出。

\begin{table}[htbp]
    \linespread{1.5}
    \zihao{5}
    \songti
    \centering
    \caption{系统功能需求表}\label{tab:functional_requirements}
    \begin{tabular}{|l|p{10cm}|}
        \hline
        功能模块            & 需求说明                                            \\ \hline
        钱包登录与用户资料管理  & 支持用户通过钱包签名完成SIWE登录，维护本地用户资料与钱包身份的一致性         \\ \hline
        创作者工作台          & 支持发行草稿创建、编辑、保存、删除和状态展示                          \\ \hline
        素材上传与元数据管理 & 支持音频、封面和元数据上传到Pinata/IPFS，并保存CID和网关地址 \\ \hline
        链上发布与分账         & 支持调用智能合约发布作品，设置访问模式、价格和收益分账地址                   \\ \hline
        链上资源检索          & 支持按标题、歌手、专辑、流派、Token ID、合约地址、交易哈希和IPFS CID检索资源  \\ \hline
        购买与访问验证         & 支持查询作品销售配置、购买访问权、验证ERC-1155访问凭证和记录购买交易          \\ \hline
        播放与链上音乐库        & 支持播放公开或已购买作品，将作品加入链上音乐库，并进行本地缓存和下载              \\ \hline
        评论互动            & 支持登录用户查看和发布作品评论                                 \\ \hline
    \end{tabular}
\end{table}


\section{非功能需求分析}

\begin{table}[htbp]
    \linespread{1.5}
    \zihao{5}
    \songti
    \centering
    \caption{系统非功能需求表}\label{tab:non_functional_requirements}
    \begin{tabular}{|l|p{10cm}|}
        \hline
        需求类型 & 说明                                                 \\ \hline
        安全性  & SIWE登录需要防止nonce重放，后端受保护接口需要进行会话校验，付费作品访问状态以链上凭证为依据 \\ \hline
        可追溯性 & 作品应通过chainId、合约地址、tokenId、CID和交易哈希进行定位和追踪          \\ \hline
        可用性  & 创作者发布流程应提供明确阶段状态和错误提示，用户购买和播放流程应尽量减少操作步骤           \\ \hline
        可维护性 & 桌面端、后端和合约应按模块划分职责，降低后续扩展成本                         \\ \hline
        性能   & 检索排序应限制候选规模，音频播放应支持Range请求和缓存机制                    \\ \hline
    \end{tabular}
\end{table}


\section{可行性分析}

\subsection{技术可行性}

Electron、React、Node.js、Express、Prisma、Solidity、ethers.js和IPFS/Pinata均为成熟技术，能够满足桌面端、后端、合约交互和资源存储需求。系统原型已经完成主要功能模块实现，因此技术上具有可行性。

\subsection{经济可行性}

系统主要基于开源技术栈实现，开发成本主要来自开发时间、测试链gas费用和Pinata存储服务。论文实现与测试阶段可使用Sepolia测试网和有限规模测试数据，经济成本可控。

\subsection{时间可行性}

系统原型已经具备主要功能，后续工作主要集中在测试完善、图表补充和论文整理。因此，在既定论文周期内完成系统说明、实验验证和论文撰写是可行的。


\section{本章小结}

本章从应用场景、用户角色、功能需求、非功能需求和可行性等方面分析了FreeFlow系统建设需求。下一章将在需求基础上给出系统总体设计。
